\documentclass[11pt,letterpaper]{article}
\usepackage{cornell-notes}
\usepackage{mathtools}

\title{Chapter 17: Strings and Evolutionary Trees}
\author{}
\date{}

\setDocTitle{Chapter 17: Strings and Evolutionary Trees}
\setDocSubtitle{String Algorithms}
\setDocVersion{v1.0}
\setDocStatus{Study Companion}
\setDocOwner{String Algorithms Cornell Notes}
\setDocAuthor{}
\setDocDate{\today}
\setDocSummary{Cornell-format study and review notes for Chapter 17: Strings and Evolutionary Trees.}

\setCornellCollection{String Algorithms}
\setCornellUnitType{Chapter}
\setCornellUnitNumber{17}
\setCornellUnitTitle{Strings and Evolutionary Trees}
\setCornellCompanionLabel{Study and review companion}

\hypersetup{pdftitle={Chapter 17: Strings and Evolutionary Trees},pdfsubject={Cornell notes},pdfauthor={}}

\begin{document}
\maketitle

\begin{CornellOverviewBox}{Chapter overview}
\textbf{Purpose.} Connect pairwise sequence distances and character changes to ultrametric, additive, parsimony, and perfect-phylogeny tree reconstruction.

\textbf{Learning objectives}
\begin{itemize}
  \item Recognize ultrametric and additive distance conditions.
  \item Reconstruct or validate a weighted tree from exact distances.
  \item Compute a parsimony score on a fixed tree.
  \item Explain the mutual dependence of alignment and phylogeny.
\end{itemize}

\textbf{Prerequisites.} Weighted trees, metrics, dynamic programming, sequence alignment, and basic evolutionary concepts.
\end{CornellOverviewBox}

\begin{CornellNotesTable}
\CornellNoteRow{What are distance- and character-based methods?}{Distance methods compress each sequence pair into one dissimilarity value and fit a tree to the matrix. Character methods retain site states and score changes along tree edges. Compression improves efficiency but can discard which sites support each branch.}
\CornellNoteRow{\S17.1: What is an ultrametric?}{An ultrametric satisfies $d(x,z)\le\max\{d(x,y),d(y,z)\}$ for every triple. Equivalently, the two largest distances in a triple are equal. It corresponds to a rooted tree whose leaves are equally distant in time from the root under a strict clock model.}
\CornellNoteRow{How is an ultrametric tree reconstructed?}{Repeatedly cluster the closest compatible taxa or distance classes and assign an ancestor height consistent with their distance. After joining a cluster, update its distance to other clusters according to the ultrametric rule. Exact input must satisfy the triple condition; noisy data require an approximation method.}
\CornellNoteRow{What is the correctness cue for clustering?}{In an exact ultrametric, a closest pair forms a valid lowest cluster because no third leaf can be closer to one member without violating the triple condition. Contracting that cluster preserves ultrametricity on the reduced instance.}
\CornellNoteRow{\S17.2: What is an additive distance?}{A distance matrix is additive if there is a weighted unrooted tree whose path lengths equal all pairwise distances. The four-point condition states that, for any four leaves, the two largest of the three pair-sums are equal. Unlike ultrametrics, leaves need not be equidistant from a root.}
\CornellNoteRow{How is an additive tree recovered?}{Identify cherries or infer limb lengths from distance differences, remove or contract a leaf, recursively reconstruct the smaller tree, and attach the leaf at the uniquely implied point. Nonnegative edge lengths and consistency checks are essential.}
\CornellNoteRow{\S17.3: What is parsimony?}{Assign states to internal nodes so the total number or cost of changes over edges is minimized. On a fixed tree, a bottom-up DP stores, for every node and possible character, the best cost for its subtree; a top-down pass recovers ancestral assignments.}
\CornellNoteRow{What is the fixed-tree parsimony recurrence?}{For node $v$ and state $a$, sum over children $u$ the minimum of $F_u(b)+c(a,b)$ over child states $b$. Leaves permit only their observed state. Columns can be scored independently when the model has no interaction across sites.}
\CornellNoteRow{Why is choosing the tree harder?}{The fixed-tree DP is polynomial, but searching tree topologies grows superexponentially and common optimization variants are computationally hard. Branch-and-bound, heuristics, and consensus across good trees are therefore common.}
\CornellNoteRow{\S17.4: Why is the ultrametric problem central?}{Ultrametric reconstruction offers a clean example in which local distance conditions completely characterize a rooted tree. It clarifies how assumptions such as a molecular clock create algorithmic tractability and how violations signal model mismatch or noise.}
\CornellNoteRow{\S17.5: How are parsimony and Steiner trees related?}{Treat possible ancestral states as Steiner vertices connecting observed leaves in a metric state space. Minimizing evolutionary change resembles a Steiner-tree problem. Perfect phylogeny is the special case where each character state arises without repeated independent changes under its model.}
\CornellNoteRow{What is perfect phylogeny testing?}{Encode each character's state partition and ask whether one tree realizes all characters without forbidden recurrence. Compatibility conditions vary with binary, multistate, ordered, or unrooted models; the character model must be stated before applying a test.}
\CornellNoteRow{\S17.6: What is phylogenetic alignment?}{Instead of treating the alignment as fixed, score insertions, deletions, and substitutions along a tree and seek ancestral sequences plus alignments on edges. This models evolution more directly but combines already difficult alignment and topology choices.}
\CornellNoteRow{\S17.7: Why do alignment and tree construction interact?}{An alignment defines the columns used to infer a tree, while a tree guides which residues should be aligned as homologous. Alternating improvement can stabilize at a local optimum, so robustness should be checked across guide trees, scoring choices, and data subsets.}
\CornellNoteRow{\S17.8: What should practice emphasize?}{Test triple and four-point conditions, reconstruct a small exact tree, run fixed-tree parsimony by hand, and identify which evolutionary assumptions are embedded in each objective.}
\end{CornellNotesTable}

\begin{CornellSummaryBox}{Chapter synthesis}
Evolutionary reconstruction can start from distances or from characters. Ultrametrics encode rooted clocklike trees through a triple condition; additive metrics encode weighted unrooted trees through a four-point condition; parsimony assigns internal character states by tree DP. When topology and alignment are also unknown, the problems become mutually dependent and computationally harder.
\end{CornellSummaryBox}

\begin{CornellSummaryBox}{Key takeaways}
\begin{CornellRecallList}
  \item An ultrametric's two largest distances in every triple are equal.
  \item An additive tree metric satisfies the four-point condition.
  \item Exact distance conditions permit recursive reconstruction.
  \item Parsimony on a fixed tree is a bottom-up state DP.
  \item Choosing a topology is much harder than scoring one.
  \item Perfect phylogeny depends on the character-state model.
  \item Alignment and phylogeny influence one another.
\end{CornellRecallList}
\end{CornellSummaryBox}

\begin{CornellOverviewBox}{Algorithm selection: when to use this}
\begin{CornellChecklist}
  \item Use ultrametric methods when a rooted clock model is credible.
  \item Use additive-distance reconstruction for exact or well-fitted path metrics without equal root-to-leaf times.
  \item Use fixed-tree parsimony to infer economical ancestral states.
  \item Use joint or iterative methods cautiously when alignment and topology are both uncertain.
\end{CornellChecklist}
\end{CornellOverviewBox}

\begin{CornellExamBox}{Self-test questions}
\begin{CornellRecallList}
  \item State and interpret the ultrametric triple condition.
  \item State the additive four-point condition.
  \item Why can the closest pair be contracted in an exact ultrametric?
  \item Write the fixed-tree parsimony recurrence.
  \item How is perfect phylogeny different from ordinary minimum parsimony?
  \item Why can changing an alignment change the inferred tree?
\end{CornellRecallList}
\end{CornellExamBox}

{\footnotesize
\textbf{Terms and notation to remember.} ultrametric, molecular clock, additive metric, four-point condition, cherry, limb length, parsimony, ancestral state, Steiner tree, perfect phylogeny.

\textbf{Connections.} Pairwise and multiple alignments supply the distances or characters used here; genome rearrangement distances broaden evolutionary comparison beyond point edits.
}

\end{document}
